% !TEX root = ../notes_sovereignai.tex
% Chapter 4: Screens and Voices. The two faces of the served model.
% Open WebUI (moved here from the stack chapter, where it was a layer it never
% was) and the Cloudflare family door (moved here from the reach chapter, which
% has no page worth publishing), then the voice pipeline: HA + Wyoming +
% Whisper + Piper + Ollama.
%
% The ordering principle: every section in this chapter is a CLIENT of the one
% Ollama. None of them holds weights, none of them touches the card, and any of
% them can be removed without the model noticing.
%%%%%%%%%%%%%%%%%%%%%%%%%%%%%%%%%%%%%%%%%%%%%%%%%%%%%%%%%%%%%%%%%%%%%%

\index{Open WebUI}\index{Cloudflare}\index{Home Assistant}\index{Wyoming}\index{Whisper}%
\index{Piper}\index{voice assistant}\index{client}
\chapter{Screens and Voices}
\printchapterversion{04_screensvoices}
\newthought{Type to It, Talk to It.}
\vspace{1.5em}

% ==============================================================================================
% THE WHY
% ==============================================================================================

\section{The Why}\index{smart home}\index{privacy}

\newthought{Ollama reaches every device you own} and nobody in the house can use
it. Reaching it means a terminal, two environment variables, and a curl, which is a
fine interface for a program and no interface at all for the person in the kitchen.
This chapter gives it the two faces a household actually uses: a page anyone
can open in a browser, and a voice that answers out loud in the room.

\marginnote{\newthought{Everything here is a client.} No section in this chapter
holds weights, touches the card, or changes what the model can do. Each one is a
way in, and any of them can be deleted without the model noticing, which is why
they were kept out of the build.}

\newthought{The screen is Open WebUI}, a browser chat with accounts and history that
points at Ollama and lists the models it already has. It is what turns one
served API into something a family uses, and it is also what publishes to a relative
who will never install a VPN.

\newthought{The voice is the harder half}, and it is the one piece of household AI
most people already rent in its worst form: an always-listening microphone that
ships every word to a company's servers. The sovereign version hears, thinks, and
speaks entirely on your hardware, and it can act on the house rather than only
answer trivia.

\newthought{The orchestrator is Home Assistant.}\index{Home Assistant!container} It is the open hub the smart-home
world already speaks, it runs as a container on the rig, and its voice pipeline is
built from swappable parts joined by a small protocol. We give it a local ear, a
local voice, and the local model as its brain.

\newthought{By the end of this chapter} you will have:
\begin{itemize}
    \item opened a browser chat on the rig and given each person in the house an account;
    \item published that chat at a hostname a relative can open with a login and no VPN;
    \item run Home Assistant with local speech-to-text and text-to-speech;
    \item made the Ollama model its conversation brain, able to act on devices;
    \item spoken to a satellite and heard the rig answer, with nothing leaving the house.
\end{itemize}

% ==============================================================================================
% THE CHAT
% ==============================================================================================

\section{The Chat}\index{Open WebUI}

\newthought{Open WebUI in a container}, a browser chat with accounts, history,
and model selection, pointed at Ollama\cite{openwebuiRepo}:

\begin{itemize}[topsep=0.25em, itemsep=0.4em, parsep=0em]
    \item For the household, it holds one account per person, so history and
    settings stay separate and nobody shares a session.
    \item For persistence, a named Docker volume\index{Docker!volume} carries accounts and
    conversations across restarts and image upgrades.
    \item For the models, it takes one variable, the base URL, and lists what
    Ollama already has. There is no second model library to maintain.
\end{itemize}

\marginnote{\newthought{The first account is admin.}\index{Open WebUI!accounts} The first sign
up becomes the administrator, so create it before anyone else opens the page,
then add an account per person from the admin panel. That same account list
mints per-person API keys\index{Open WebUI!API keys}, which is how you hand one member of the household
programmatic access without enrolling their machine on your mesh.}

\newthought{Run it and sign in.} The container reaches Ollama on the host
through \texttt{host.docker.internal}, which the added host maps to the gateway
address:

\begin{verbatim}
$ docker run -d --name open-webui -p 3000:8080 \
    --add-host=host.docker.internal:host-gateway \
    -e OLLAMA_BASE_URL=http://host.docker.internal:11434 \
    -v open-webui:/app/backend/data \
    ghcr.io/open-webui/open-webui:main
\end{verbatim}

\marginnote{\newthought{It is already everywhere.}\index{tailnet} A device on your mesh reaches
any port on the rig, not only the one served over HTTPS, so \texttt{rig:3000}
opens from a phone on cellular with nothing further to configure. The section
below is for the people who are not on the mesh.}

\newthought{What you now own} is a chat anyone in the house can open. From any
browser on the home network, go to \texttt{http://rig-address:3000}, sign in, and
ask it something. A family member two rooms away is talking to a model running on
your 3090, with no token metered and no prompt leaving the house.

% ==============================================================================================
% THE FAMILY DOOR
% ==============================================================================================

\section{The Family Door}\index{Cloudflare!Tunnel}

\newthought{A Cloudflare tunnel publishes the chat} to a relative who wants a link
and will never install a VPN\cite{cloudflareTunnel}:

\begin{itemize}[topsep=0.25em, itemsep=0.4em, parsep=0em]
    \item For reach, a small daemon on the rig dials out to Cloudflare, so the
    hostname works with no port forwarded and no inbound rule on your router.
    \item For transport, Cloudflare issues and terminates the TLS\index{TLS termination}, so the link is
    a real \texttt{https} address rather than a certificate warning.
    \item For the door itself, Cloudflare Access\index{Cloudflare!Access} puts a login in front and takes
    a list of email addresses, so the page is public and the chat is not.
\end{itemize}

\marginnote{\newthought{The sovereignty cost is real.} Cloudflare terminates the TLS
and sees who connects and when, though never what runs on your hardware. That is
the trade for a link a relative can open, and it is the only piece of this book
that puts a company back in the path. Publish the chat and nothing else.}

\newthought{Publish one hostname} and point it at the chat container, never at
Ollama:

\begin{verbatim}
$ cloudflared tunnel login
$ cloudflared tunnel create rig
$ cloudflared tunnel route dns rig chat.example.com
$ cloudflared tunnel run --url http://localhost:3000 rig
\end{verbatim}

\newthought{Add the policy before you share the link.} In the Cloudflare dashboard,
under Zero Trust, Access, Applications, add the hostname and an allow policy
limiting it to your family's email addresses. Until that policy exists the page is
open to anyone who guesses the name.

% ==============================================================================================
% SPEECH TO MODEL TO SPEECH
% ==============================================================================================

\section{Speech to Model to Speech}\index{speech to text}\index{text to speech}%
\index{STT|see{speech to text}}\index{TTS|see{text to speech}}

\newthought{A voice turn is three steps.} Speech-to-text turns the spoken request
into words; the model decides what it means and what to do; text-to-speech turns
the answer back into sound. The local stack is faster-whisper for the ear, the
model already served on the rig for the brain, and Piper for the voice, each a
small service, each replaceable.

\marginnote{\newthought{Wyoming}\index{Wyoming!protocol} is Home Assistant's protocol for voice parts: each
service listens on a TCP port and speaks audio and events, and Home Assistant wires
them into a pipeline. Whisper listens on \texttt{10300}, Piper on \texttt{10200}.}

\newthought{Put the ear on the GPU.}\index{faster-whisper} Speech-to-text is the slow step, and the
built-in add-on runs on the CPU; running faster-whisper as its own container with
the GPU passed through brings a request under a second. The same 3090 that serves
chat transcribes speech in the gaps, since the two rarely peak at once in a house.

\marginnote{\newthought{The brain must hold tools.}\index{tool calling}\index{conversation agent} A conversation agent that only
chats can answer the weather but cannot turn off a light. Home Assistant exposes its
devices to the model as tools, so the model must be one that supports tool calling.
The Qwen model on the rig was pulled for exactly this.}

% ==============================================================================================
% WIRE THE PIPELINE
% ==============================================================================================

\section{Wire the Pipeline}\index{voice pipeline}

\newthought{Run the orchestrator and the voice parts} as containers on the rig:
Home Assistant, the ear with the GPU passed through, and the voice.

\begin{verbatim}
$ docker run -d --name homeassistant --network host \
    -v ha-config:/config ghcr.io/home-assistant/home-assistant:stable

$ docker run -d --name whisper -p 10300:10300 --gpus all \
    rhasspy/wyoming-whisper --model small-int8 --language en --device cuda

$ docker run -d --name piper -p 10200:10200 \
    rhasspy/wyoming-piper --voice en_US-lessac-medium
\end{verbatim}

\marginnote{\newthought{Add the parts in the UI.}\index{Home Assistant!Wyoming integration} In Home Assistant, Settings,
Devices and Services, add the Wyoming Protocol integration twice, pointing it
at the rig's address on ports \texttt{10300} and \texttt{10200} to register the ear
and the voice.}

\newthought{Make the model the brain.}\index{Home Assistant!Ollama integration} Add the Ollama integration and point it at
the rig, then choose a tool-calling model:

\begin{verbatim}
Settings > Devices & Services > Add Integration > Ollama
  URL ..... http://localhost:11434
  Model ... qwen3.5:9b
  [x] let the model control Home Assistant devices
\end{verbatim}

\newthought{Assemble the pipeline.}\index{Home Assistant!voice pipeline} Under Settings, Voice Assistants, build one
assistant from the three local parts: faster-whisper for speech-to-text, the Ollama
agent for conversation, and Piper for text-to-speech. Set it as the preferred
assistant.

\marginnote{\newthought{Give it ears in a room.}\index{voice satellite}\index{Wyoming!satellite} A Home Assistant Voice PE puck, or
any Wyoming satellite, is the microphone and speaker on the shelf; it hands audio to
the pipeline and plays the answer back. One per room you want to talk to.}

\newthought{What you now own} answers out loud. You say ``turn off the kitchen light
and tell me a joke.''
The satellite streams it to the rig, faster-whisper transcribes it, the model turns
off the light through Home Assistant and writes a reply, and Piper says it back, all
on hardware you own, with the router unplugged if you like.

% ==============================================================================================
% A MODEL PER FACE
% ==============================================================================================
% The in-depth model selection sits HERE and not in the stack chapter. The stack
% chapter gives the sizing law and a shortlist for one card, which is all a reader
% needs to get a first answer. Choosing per task only becomes a real decision once
% the tasks exist, and both faces exist by this point in the chapter.

\section{A Model per Face}\index{model selection}\index{Ollama!concurrency}

\newthought{The two faces want different models.} The browser chat is read at
reading speed and rewards the largest model that fits, because nobody minds
waiting two seconds for a better answer. The voice pipeline is judged on the gap
between the end of your sentence and the start of the reply, and a model that
answers better but starts later is the worse assistant in a kitchen.

\marginnote{\newthought{Latency is the voice budget.} A spoken turn already spends
time on transcription and on speech synthesis. Whatever the model adds lands on
top of that, and past roughly two seconds the person in the room repeats
themselves, which starts the turn over.}

\newthought{Three roles, three shapes.}\index{embedding model} A chat model, a tool caller, and an
embedding model are different jobs, and only the first is a matter of taste. A
tool caller has to have been trained to emit structured calls, and an embedding
model does not converse at all: it turns text into vectors and is measured on
dimensions and on how fast it can chew through a folder.

\vspace{0.5em}
\noindent{\setlength{\tabcolsep}{4pt}
\begin{tabular}{p{0.13\textwidth}p{0.35\textwidth}p{0.12\textwidth}p{0.27\textwidth}}
    \toprule
    \textbf{Face} & \textbf{Model} & \textbf{Size} & \textbf{Bound by} \\
    \midrule
    Chat & \texttt{qwen3.8:27b} & 18 GB & answer quality \\
    Voice & \texttt{qwen3.5:9b} & 6.6 GB & latency \\
    Vision & \texttt{qwen3-vl:8b} & 6.1 GB & image support \\
    Notes & \texttt{qwen3-embedding:0.6b} & 639 MB & throughput \\
    \bottomrule
\end{tabular}}
\vspace{0.5em}

\newthought{They share one card, up to a point.}\index{VRAM} Ollama keeps several
models resident at once as long as they fit, so the embedder sits beside a chat
model without either being unloaded. What does not fit is two large ones:
\texttt{qwen3.8:27b} and \texttt{qwen3.5:9b} together want 24.6~GB of a 24~GB card,
so every switch between the browser and the satellite costs a reload, and that
reload is the pause the person in the kitchen hears\cite{ollamaLibrary}.

\marginnote{\newthought{The cheap way out} is one model for both faces.
\texttt{qwen3.5:9b} answers the browser well enough and calls tools quickly
enough, and at 6.6~GB it leaves room for the vision model and the embedder
resident beside it. Split the faces only after a task has failed on the shared
model.}

\newthought{Set the three variables} that decide how the card is shared, then let
Ollama place the models\cite{ollamaDocs}:

\begin{verbatim}
$ sudo systemctl edit ollama
\end{verbatim}

\begin{verbatim}
[Service]
Environment="OLLAMA_MAX_LOADED_MODELS=3"
Environment="OLLAMA_NUM_PARALLEL=2"
Environment="OLLAMA_KEEP_ALIVE=30m"
\end{verbatim}

\marginnote{\newthought{What each one costs.} Loaded models default to three per
card and only load if they fit. Parallel requests default to one, and raising it
multiplies the context memory by that number. Keep alive defaults to five
minutes, which is short enough that the first spoken request of the evening waits
through a load.}

\newthought{Confirm what is resident} before blaming a model for being slow. If
the tag you expect is missing, the pause you heard was a load and not the model
thinking:

\begin{verbatim}
$ ollama ps
$ ollama pull qwen3-embedding:0.6b
\end{verbatim}

% ==============================================================================================
% BRIDGE
% ==============================================================================================

\newthought{The house can now listen and speak}, and browser, terminal, and voice
all draw on the one Ollama. What it still lacks is your own life: it answers
from what it was trained on, never from what you have written. The next chapter
points your notes at it.

\marginnote{\newthought{feedback}}
